\documentclass[14pt]{extarticle}
\usepackage{preamble}
\geometry{letterpaper, landscape, margin=0.5in}
\usepackage{qrcode}
\renewcommand{\answer}[1]{}
\setlength{\parskip}{4pt}

\begin{document}
\pagestyle{empty}

\daybanner{Topic 5.5 \textperiodcentered\ TODAY'S PROBLEM}{Wrong Shape Does Not Mean Biased}

\noindent\begin{minipage}[t]{0.57\linewidth}
\vspace{0pt}
Suppose a video platform has 5,000 active accounts and says that $p=0.08$ of them regularly use captions. An SRS of 100 accounts, selected without replacement, contains 14 regular caption users, so $\widehat p=0.14$.\par\smallskip \textbf{Jules:} ``Because $n=100>30$, $\widehat p$ is approximately normal. I get $P(\widehat p\geq0.14)=0.0135$, which proves caption use increased.''\par \textbf{Nora:} ``Because $np=8<10$, a normal model is not justified. Therefore $\widehat p$ is biased and this sample tells us nothing.''\par
\end{minipage}\hfill
\begin{minipage}[t]{0.40\linewidth}
\vspace{0pt}
\begin{center}
\textbf{Caption-use sample}\par\smallskip
\begin{tabular}{|l|c|c|c|}
\hline
\textbf{Source} & \textbf{Accounts} & \textbf{Users} & \textbf{Rate} \\ \hline
\textbf{Population} & 5,000 & 400 & 0.08 \\ \hline
\textbf{SRS} & 100 & 14 & 0.14 \\ \hline
\end{tabular}
\end{center}
\end{minipage}

\begin{firsttakebox}{FIRST TAKE (5 min, on your sheet)}
Whose reasoning is closer to correct, Jules's or Nora's? Name one condition or claim you would check.
\end{firsttakebox}

\begin{description}[leftmargin=1.15in, labelwidth=1in, itemsep=2pt, topsep=2pt]
  \item[\dokbadge{1}~(a)] Identify the population parameter, sample statistic, population size, and sample size. Verify the reported value of $\widehat p$.
  \item[\dokbadge{2}~(b)] Find the mean and standard deviation of the sampling distribution of $\widehat p$. Check the 10 percent condition and both Large Counts values.
  \item[\dokbadge{3}~(c)] Adjudicate both claims. Use a binomial calculation or valid simulation for $P(\widehat p\geq0.14)$ under $p=0.08$. State the strongest conclusion warranted, explain why failed normality does not create bias, and give the consequence of each overclaim.
\end{description}

\vfill
\noindent\begin{minipage}{0.84\linewidth}
\textbf{Video + follow-along:} \href{https://robjohncolson.github.io/apstats-live-worksheet/u5_lesson5_live.html}{\texttt{\detokenize{u5_lesson5_live.html}}} (scan the code)\par
\textbf{Finish (a)--(c) after the video. Turn in your sheet.}
\end{minipage}\hfill
\qrcode[height=0.75in]{https://robjohncolson.github.io/apstats-live-worksheet/u5_lesson5_live.html}

\end{document}
